A principal dificuldade para treinar um classificador de dígitos de hidrômetro é a
rotulagem: os conjuntos públicos informam a leitura completa do medidor, e não o valor de
cada rolete. Três fontes de dados foram combinadas, cada uma com um tratamento próprio
para gerar recortes rotulados de um único dígito:

\begin{enumerate}
    \item \textbf{Water Meters Dataset} \cite{kucev2020}: 1.244 fotos de hidrômetros
    variados, tiradas em ângulos arbitrários, com a máscara da janela de dígitos e a
    leitura no nome do arquivo. Para o classificador, passou por um processo automático: a janela foi retificada a partir da
    máscara e dividida em células por projeção vertical. Para evitar rótulos errados, só
    foram aceitas as fotos em que a divisão produziu exatamente oito células de larguras
    semelhantes; as demais foram descartadas. Restaram 195 fotos, que geraram 1.560 recortes.
    \item \textbf{Word-Wheel Water Meter Dataset} \cite{zhao2026wordwheel}: conjunto público
    cujo subconjunto de reconhecimento traz recortes padronizados de $200 \times 60$ pixels
    da janela de dígitos, com a leitura completa. Foram descartadas as imagens com um rolete
    no meio da troca entre dois dígitos, cujo rótulo é ambíguo, e as imagens espelhadas
    foram invertidas. Como os recortes já vêm retos, cada um foi dividido em fatias de
    largura igual. Apenas as pastas de treino foram usadas, gerando 47.478 recortes; as
    pastas de teste, com 1.894 imagens, foram reservadas como teste independente.
    \item \textbf{Fotos do próprio \gls{esp32cam}}: para aproximar o modelo das condições
    do sensor, da lente e da compressão do dispositivo, três fotos capturadas pelo
    \gls{esp32cam}, de dois medidores diferentes, tiveram a janela marcada manualmente e
    foram fatiadas e rotuladas à mão, gerando 24 recortes. Esses recortes foram usados
    apenas no ajuste fino (\textit{fine-tuning}) do classificador.
\end{enumerate}